\section{Core Architectural Design and Execution Flow}
\label{sec:Architecture}
\vspace{2em}
\subsection{Architectural Design Patterns for Modularity}
The codebase's central logic is structured around established software design patterns to ensure high modularity, extensibility, and maintainability across heterogeneous components.

\paragraph{\textbf{Adapter Pattern} (\texttt{llm\_wrappers.py}):} This pattern standardizes the interface for interacting with heterogeneous LLM providers (e.g., Gemini, Watsonx). By abstracting the provider-specific connection details, it enables the core execution engine to interact with all models through a uniform API.

\paragraph{\textbf{Factory Pattern} (\texttt{llm\_factory.py}):} Implemented as a simple factory method, this pattern decouples client components (e.g., \texttt{sql\_baseline.py}, \texttt{nl\_baseline.py}) from the specific instantiation logic of concrete LLM wrapper objects. Extensibility to new providers is achieved solely by updating the internal \texttt{PROVIDER\_MAP} registry.

\paragraph{\textbf{Template Method Pattern} (\texttt{LLMBaseWrapper} in \texttt{llm\_wrappers.py}):} The abstract \texttt{LLMBaseWrapper} defines the invariant sequence (the *template*) for the life cycle of the LLM instance retrieval. It establishes the mandatory sequence while deferring the implementation of the provider-specific connection details (via \texttt{create\_llm\_instance()}) to concrete subclasses.

\paragraph{\textbf{Chain/Pipeline Pattern (LCEL)} (\texttt{build\_lcel\_chain.py}):} This pattern, realized using LangChain Expression Language (LCEL), defines a structured, sequential execution flow:
\begin{center}
    \textbf{Context Generation} $\to$ \textbf{Prompt Construction} $\to$ \textbf{LLM Invocation}
\end{center}
This structured approach ensures predictable and auditable data transformation across the execution path.
\newpage

\subsection{Unified Architecture and Execution Workflow (The Pipeline)}
\label{sec:core_arch}
To ensure high modularity, extensibility, and maintainability, we adopted a unified base architecture that serves both baselines (NL and SQL). This architecture centralizes the execution logic, abstracting interactions with LLM providers and standardizing the data flow.
\\\\
The execution workflow is structured into a four-stage pipeline, ensuring robust and consistent processing across all baseline types:

\begin{enumerate}[leftmargin=*,label=\arabic*.]
    \item \textbf{Context Information Retrieval}: The \texttt{get\_db\_context} function dynamically retrieves the necessary database context that will be attached to the LLM prompt together with the corresponding SQL query fro SQL baseline or NL prompt for NL baseline. For instance this includes retrieving the database schema for SQL/NL baselines.
    \item \textbf{Prompt Construction and Injection}: The derived context is programmatically injected into a sophisticated \texttt{ChatPromptTemplate}. This stage manages the insertion of system instructions and the baseline-specific input (the NL or SQL query), and also the enforce the required JSON output structure via an \texttt{PydanticOutputParser} object.
    \item \textbf{LLM Invocation}: The fully prepared and structured prompt is transmitted to the configured LLM instance, which was instantiated via the Factory/Adapter layers.
    \item \textbf{Response Standardization}: The \texttt{parse\_llm\_response} function enforces cross-provider output standardization: This involves the mandatory validation and parsing of the requisite JSON structure returned by the LLM. It then enriches this output into the \texttt{FULL\_JSON} format by extracting the \texttt{token consumption} from the LLM's response metadata and internally calculating and appending the \texttt{LLM response latency}.
\end{enumerate}
This last standardized and enriched data process is essential for downstream performance evaluation, furthermore thanks to this centralized architecture, the execution logic is provider-independent, and the workflow is uniform.

\vspace{2em}
\subsection{Data Context and Output Standardization}

This section defines the two primary contexts used to condition the Large Language Models (LLMs) and establishes the \textbf{standardized JSON output format} mandatory for all baseline evaluations (NL, SQL, and Palimpzest).

\medskip

\subsubsection{Data Context Categorization}

The evaluation utilizes two distinct dataset contexts, governing the type of factual information provided to the LLM during prompting:

\begin{itemize} [leftmargin=*]
    \item \textbf{Internal Knowledge (IK):} In these datasets, the model relies exclusively on its pre-trained knowledge and linguistic reasoning capabilities. The prompt includes only the natural language question and the database schema for both the \textbf{SQL} and \textbf{NL} baselines, but no factual data retrieved from the database.
    \item \textbf{Model Context (MC):} These datasets are designed to enrich the prompt with external information retrieved from a knowledge base. For this specific project task, MC datasets were \textbf{not used} in either the NL or SQL baselines as indicated in GALOIS paper, but could represents a good resource for RAG purposes.
\end{itemize}

\medskip

\subsubsection{Standardized Output Format}

Every response produced by any baseline is written in a JSON file that fits the mandatory \textbf{FULL JSON} format. This rigorous structure ensures consistent computation of metrics across all providers and baselines.

\begin{figure}[h!]
    \centering
    \begin{verbatim}
    {
        "result_set": [
            { "originaltitle": "The Three Musketeers"},
            { "originaltitle": "The Count of Monte Cristo"}
        ],
        "time": 1.84,
        "tokens": 312
    }
    \end{verbatim}
    \caption{\texttt{FULL JSON} Output Format Example}
\end{figure}

\begin{table}[h]
    \centering
    \begin{tabular}{l p{0.7\textwidth}} 
        \toprule
        \textbf{Field Name} & \textbf{Description} \\
        \midrule
        \textbf{result\_set} & Denotes the structured LLM output. This must be an array of JSON objects corresponding to the query's selected columns. \\
        \textbf{time} & Represents the total time taken by the LLM to generate the response. \\
        \textbf{tokens} & Indicates the number of tokens processed during the interaction. \\
        \bottomrule
    \end{tabular}
    \captionof{table}{Definition of the \texttt{FULL JSON} Output Schema Fields}
\end{table}
